\chapter{Algorithms as Inference Rules}\label{app:infer}

There are a number of ways to express algorithms in an article of computer science.
Some traditional way writes imperative pseudo code that emphasizes state changing (assignments) and loops, as shown in \autoref{algo:match-dfs-cons}.
While pseudo code in this style can be straightforwardly translated to a common imperative programming language, such as {\CPP},
it is less mathematically formal and rigorous, with a big burden to reason about.
On the other hand, functional style algorithms use equations to introduce new variables and bindings, without changing the value bound to a variable.
It is obvious that updating the state of a variable is only useful when running in a loop.
The functional style employs recursive calls to achieve much more powerful repetitions and bind new values to variables in inner calls.
It is a more concise and rigorous way to express algorithms, and a better form to reason about the correctness and properties.

\paragraph{Inference rules}
Functional style algorithms often use pattern matching with guards to perform case analyses on the input arguments, and
bind new variables to components and intermediate results for further specification.
Just like the chained conditional statements in traditional pseudo code, these cases are listed alongside the left for clarity.
This is space wasting for publishing articles in journals and proceedings.
An elegant space saving way of listing the possibly nested cases is to formulate them as a sequence of {\em inference rules},
\begin{mathpar}
\inferrule[(Case 1)]
  {A \\ B \\ C}
  {\inferrulel[(Inner Case)]
    {F \\ G}
    {K_1} \\
   \inferrulel[(Inner Else)]
    {\quad}
    {K_2}
  }
\and
\inferrule[(Case 2)]
  {\quad D \\ E \quad}
  {K_3}
\and
\inferrule[(Else)]
  {\quad \\ \quad}
  {K_4},
\end{mathpar}
where $A,B,C,D,E$ are the top-level premises, $F,G$ the inner-level premises of the rule labeled ({\sc Inner Case}), $K_1, K_2, K_3, K_4$ the conclusions.
Inference rules are examined sequentially, thus, $\neg(A \wedge B \wedge C)$ is implied in ({\sc Case 2}).
For example, the Euclid's greatest common divisor algorithm $\mathit{gcd}(m, n)$ can be defined on integers $m,n$ as follows,
\begin{mathpar}
\mathit{gcd}(m,n) =
\and
\inferrule
  {m < 0}
  {\mathit{gcd}(-m,n)}
\and
\inferrule
  {n < 0}
  {\mathit{gcd}(m,-n)}
\and
\inferrule
  {m < n}
  {\mathit{gcd}(n, m)}
\and
\inferrule
  {n = 0}
  {m}
\and
\inferrule
  {\quad}
  {\mathit{gcd}(n, m-n)}.
\end{mathpar}

Generally, functional style algorithms are more abstract and less to write.
However, for accumulation over a collection of elements, some verbose arrangement is necessary to explicitly handle the iteration of elements.
Let's see another example to translate $\DfsMatch$ in \autoref{algo:match-dfs-cons} to inference rules.
\begin{mathpar}
\DfsMatch'(f, \emptyset, E, E') = f,
\\
\DfsMatch'(f, \lstt{(p, p')}\uplus P, E, E') =
\and
\inferrule
  {p \in \dom(f)}
  {\inferrule{f(p) \ne p'}{\KwRaise\ \Mismatch} \\
   \inferrule{\quad}{\DfsMatch'(f,P,E,E')}
  }
\and
\inferrule
  {\tau(p') \nsubty \tau(p)}
  {\KwRaise\ \Mismatch}
\and
\inferrule
  {p' \in \ran(f)}
  {\KwRaise\ \Mismatch}
\and
\inferrule
  {\exists p \la{a} q \in E \suchthat p' \la{a} q' \not\in E'}
  {\KwRaise\ \Mismatch}
\and
\inferrule
  {Q = \lstt{(q, q') \mid p \la{a} q \in E \wedge p' \la{a} q' \in E'}}
  {\DfsMatch'(f\cup\sett{p \mapsto p'},Q \uplus P,E,E')}.
\end{mathpar}
The use of list concatenation $\uplus$ in the above definition implements the stack pop and push that guarantee the \DFS.

\paragraph{Overlined terms}
To reduce the occasions of explicit enumeration of elements over a collection, {\em overlined terms} are often used in certain cases.
To the simplest case, an overlined varible $\overline{v}$ is expanded to a series of variables $v_1,\dots,v_n$.
This expansion is usually restricted by the context, where a component may already be defined.
A variable that is defined as a singlton repeats itself in the expansion.
A variable that is also a series pairs with the other series element-wise in the expansion.
For example, let $f$ be a defined function, $\overline{a}$ a series. Then $\overline{b = 2a}$ expands to $b_1 = 2a_1,\dots,b_n = 2a_n$, and
$\overline{f(a,b)}$ expands to $f(a_1,b_1),\dots,f(a_n,b_n)$.
With the overlined terms, $\DfsMatch'$ can be re-written without the list comprehension.
\begin{mathpar}
\DfsMatch'(f, \lstt{(p, p')}\uplus P, E, E') =
\and
\inferrulel[(visited)]
  {f(p) = p'}
  {\DfsMatch'(f,P,E,E')}
\and
\inferrulel[(explore)]
  {\tau(p') \subty \tau(p) \\ \overline{p \la{a} q \in E \implies p' \la{a} q' \in E'} \\\\ p \notin \dom(f) \\ p' \notin \ran(f)}
  {\DfsMatch'(f\cup\sett{p \mapsto p'},[\overline{(q,q')}] \uplus P,E,E')}
\and
\inferrule
  {\quad}
  {\KwRaise\ \Mismatch}.
\end{mathpar}
I would say the overline used in this definition is quite intuitive. If a detailed explanation must be given,
$p, p', E, E'$ are bound variable in the context, $a, q, q'$ are new series.
Since $q$ and $q'$ are all aligned with $a$, they have the correspondence synchronized by $a$.
Therefore, we can expand the overlined term as, for all the outgoing edge labels $a_1,\dots,a_n$,
\[
p\la{a_1}q_1 \in E \implies p'\la{a_1}q_1' \in E', \quad\dots,\quad p\la{a_n}q_n \in E \implies p'\la{a_n}q_n' \in E',
\]
and we have the list of pairs $[(q_1,q_1'), \dots, (q_n,q_n')]$ in the conclusion.

I must admit that the rules for the expansion are not rigorous. There are some ambiguities that require common sense to remove.
What I really mean in the above definition is
\[
\forall p\la{a}q \in E \suchthat p'\la{a}q' \in E',
\]
but it cannot be written as a simple series of prepositions, without specifying the range of $a$ in a verbose way.
I don't want to sacrifice the simplicity for the negligible bit of additional rigor.
At the end of the day, we can still complement the specification with remarks in the text.

A formally available construct for expandable series is the {\em variadic templates} in {\CPP}~\cite{ISO:2023:n4950}.
It provides some great ideas for expanding series and pairing variables.
An exciting feature of this construct is the {\em fold expression} for iteration and accumulation of elements in a series.

\begin{dispCPP}
template<class... T>
void print(T... a)
{
    (std::cout << ... << a) << std::endl;
    ((std::cout << ' ' << a), ...) << std::endl;
}

int main()
{
    print('x', 100, true, 99.87);
    return 0;
}
\end{dispCPP}

\endinput
